%% LyX 2.4.4 created this file.  For more info, see https://www.lyx.org/.
%% Do not edit unless you really know what you are doing.
\documentclass[english]{article}
\usepackage[T1]{fontenc}
\usepackage[latin9]{inputenc}
\usepackage{graphicx}
\usepackage{geometry}
\geometry{verbose,tmargin=1in,bmargin=1in,lmargin=1in,rmargin=1in}

\makeatletter
%%%%%%%%%%%%%%%%%%%%%%%%%%%%%% User specified LaTeX commands.
\usepackage{hyperref}
\usepackage[usenames,dvipsnames]{xcolor} %adds additional color options
\hypersetup{colorlinks=true, urlcolor=Blue, linkcolor=Blue, citecolor=Blue}
\usepackage{url}

\usepackage{xspace} %\xspace in \newcommand removes the need for {} to get proper spacing
\usepackage{multicol}
\usepackage{advdate}    % Advancing/saving dates
\usepackage{datetime}   % Dates formatting
\usepackage{datenumber} % Counters for dates
\usepackage{hyphenat} % custom hyphenation ex. $\beta$\hyp{}galactosidase

\makeatother

\usepackage{babel}
\begin{document}
\newcommand{\me}{Dr.~Ashe\xspace}
\newcommand{\meL}{Dr.~James Ashe\xspace}
\newcommand{\guest}{Dr.~Stephenson\xspace}
\newcommand{\guestL}{Dr.~Ken Stephenson\xspace}
\usdate

\begin{center}
{\large\textbf{2013 HBCU-UP Mini Grants Program }}{\large\par}
\par\end{center}

\begin{center}
{\large\textbf{Request for Proposals}}\\
\par\end{center}

\section{\label{sec:Purpose-and-specific}Purpose and specific hypotheses/research
questions}

Circle packings are configurations of circles with specified patterns
of tangency. They naturally express a conformal geometry which has
allowed the field of circle packing to be developed into a rich theory
of discrete analytic functions, see \cite{Ste:MR2131318,Ste:}. The
discretization of complex analysis via computable circles has application
advantages. For example, the NSF has supported research for using
circle packings in brain imaging, see \cite{hurdal2004cortical,hurdal1999quasi,hurdal2001case}. 

Using circle packings \guestL has constructed discrete models of
the Thomson problem. A longstanding problem in physics with applications
in biology, the Thomson problem asks how to find the minimum energy
distribution of a finite number on electrons on the surface of the
sphere. By applying a Delaunay triangulation to any distribution,
an approximate spherical circle packing model can be found, see figure
\ref{fig:An-energy-distribution}. Furthermore, combinatorial manipulations
can be used to minimize the total energy. It is hoped that exploring
these manipulations will provide insights into the problem. 

\begin{figure}
\hfill\includegraphics[height=5cm]{Admin/images/Thomson_410_mesh_thermal.png}\hspace{1cm}\includegraphics[height=5cm]{Admin/images/cpShot-001.jpg}\hfill{}

\caption{\label{fig:An-energy-distribution}An energy distribution modeled
with circle packing.}
\end{figure}


\subsection{What will be done}

We propose that Johnson C.\,Smith undergraduate students be tasked
with exploring the Thomson problem via circle packings. Students should
be able to construct discrete models and measure their approximation
of the classical problem. Through combinatorial manipulation, they
will experiment with distributions and seek to make conjectures about
the problem. 

\subsection{Estimated start-up and completion dates}

It is expected that work can begin by \monthname[2] 2013 and continue
until \monthname[5] 2013. Because of the variety of possible experiments,
there is potential that this work may develop into a larger project. 

\subsection{How}

See Section \ref{sec:Procedures-for-evaluation} below. 

\subsection{By Whom}

The necessary code for implementing and manipulating the spherical
packing models will be written by \guest and \me. Thus the student(s)
would need little prerequisite programming or mathematical knowledge,
only a willingness to learn and the patience to conduct the experiments. 

JCSU student Daniel Villanueva and Merischia Griffin have been recommended
for this project. It is an excellent opportunity to be introduced
to programming and high level mathematics.

\section{\label{sec:Procedures-for-evaluation}Procedures for evaluation}

Circle packings can be manipulated by Java\texttrademark code or
directly through ${\tt CirclePack}$. A configuration of circles (location
an radii) can then be computed with ${\tt CirclePack}$. Energy points
are identified with circle vertices, and energy is then easily calculated
with software. Although there has been little study of the Thomson
problem using circle packings, the problem has been thoroughly researched.
Thus there is an abundance of available data for the derivation of
discrete models. 

The creation and analysis of these experiments are done strictly on
computers. The following materials will be required:
\begin{itemize}
\item ${\tt CirclePack}$; creates and manipulates circle packings. This
software is open source and freely available. 
\item Eclipse; a multi-language software development environment. This software
is open source and freely available. 
\item Java\texttrademark; ${\tt CirclePack}$'s programming language. This
software is freely available. 
\item ${\tt MATLAB}$\textregistered; a numerical computing and analysis
program.
\item Suitable computer hardware for using the above software.
\end{itemize}

\section{Budget}

The required software ${\tt CirclePack}$, Java\texttrademark, and
Eclipse are freely available through public licenses. Student and
faculty licenses for ${\tt MATLAB}$\textregistered\,should be available
through existing JCSU resources. Thus costs will be limited mostly
to compensation and travel expenses. An estimate is as follows:
\begin{itemize}
\item \$1000: Stipend for \me.
\item \$1000: Stipends for two students (\$500 each).
\item \$1482: Travel allowance for students and \me. 
\item \$518: Hosting \guest at JCSU from \formatdate{10}{2}{2013} to \formatdate{12}{2}{2013}.

\begin{itemize}
\item \$238 (\$119 per night): Lodging at the Mosaic Village for \guest
\item \$175.04 (233.8 miles at \$0.375 per mile): \guest's mileage cost
from his home in Knoxville, TN to JCSU.
\item \$105 (\$35 per day): per diem for \guest. 
\end{itemize}
\item \$300: Computer hardware (this cost may be reduced or eliminated with
corporation with the IT department).

\begin{itemize}
\item An additional monitor.
\item Graphics card capable of supporting dual monitors and handling 3D
modeling.
\end{itemize}
\item \$100: Presentation materials.
\end{itemize}
Additional expenses to be paid for by HBCUP:
\begin{itemize}
\item \$1000: Stipend for visiting scientist \guest. This cost also pays
\guest's honorarium for his visit on \formatdate{10}{2}{2013} to
\formatdate{12}{2}{2013}. 
\end{itemize}

\section{Time required}

There is a very large number of possible experiments, and anything
discovered would be of value. Thus students will be compensated at
the standard rate until their stipends are exhausted. 

Due to the nature of programming and experimental mathematics, the
time required by the faculty is difficult to estimate. In addition
to programming, the faculty would need to advise students, create
images, and attach mathematical rigor to the students' experimental
results. An estimate of 40+ hours for the faculty advisors might be
conservative. Thus the stipend compensation will most certainly cost
less than that of a programmer/mathematician being paid at the standard
rate. 

\section{Ethical Issues}

There are no known ethical issues particular to this project.

\section{Benefits of Research}

Benefits to the study of mathematics are outlined in section \ref{sec:Purpose-and-specific}.
The students will simultaneously gain experience in programming, algorithms,
optimization, and pure mathematics. Results might be publishable,
garner note in a high level math publication, or be the cornerstone
of a senior project. Additionally, the work would enhance a r\'{e}sum\'{e},
graduate school application, or portfolio.

\section{Plan to Share and Disseminate}

This visual nature of circle packings are excellent for capturing
the attention of mathematicians and lay people alike. Posters and
the junior seminar would be appropriate formats to display and discuss
early results. Further along, regional undergraduate math conferences
and math journals would be the next step. \guest will visit JCSU,
work with students, and speak at a STEM College lyceum on \formatdate{11}{2}{2013}. 

\bibliographystyle{plain}
\bibliography{/home/jim/JamesAshe/bib_jashe}

\end{document}
